\chapter{Algorithms}
\label{chap:algorithms}

\section{Time Complexity}
Time complexity for an algorithm has three different facets
\begin{itemize}
    \item The \highlightit{worst-case} complexity of the algorithm is 
          the function defined by the maximum number of steps taken in 
          any instance of size $n$. 
    \item The \highlightit{best-case} complexity of the algorithm is 
          the function defined by the minimum number of steps taken in 
          any instance of size $n$. 
    \item The \highlightit{average-case} complexity or expected time of 
          the algorithm, which is the function defined by the average 
          number of steps over all instances of size $n$. 
\end{itemize}
\newpage
\subsection{Big O notation}
There exist constants $c > 0$ and $n_0$ such that
\begin{equation}
    \label{eqn:bigOnotation}
    f(n) = O(g(n)) \quad \text{if} \quad 0 \le f(n) \le c\, g(n)
    \quad \forall\, n \ge n_0 .
\end{equation}

For all values $n \ge n_0$, $f(n)$ is on or below $c\, g(n)$.
\newline
A constant multiple of $g(n)$ is an \highlightit{asymptotic upper bound} for $f(n)$.

\subsection{Omega notation}
There exist constants $c > 0$ and $n_0$ such that
\begin{equation}
    \label{eqn:omeganotation}
    f(n) = \Omega(g(n)) \quad \text{if} \quad 0 \le c\, g(n) \le f(n)
    \quad \forall\, n \ge n_0 .
\end{equation}

For all values $n \ge n_0$, $f(n)$ is on or above $c\, g(n)$.
\newline
A constant multiple of $g(n)$ is an \highlightit{asymptotic lower bound} for $f(n)$.

\subsection{Theta notation}
There exist constants $c_1, c_2 > 0$ and $n_0$ such that
\begin{equation}
    \label{eqn:thetanotation}
    f(n) = \Theta(g(n)) \quad \text{if} \quad
    0 \le c_1\, g(n) \le f(n) \le c_2\, g(n)
    \quad \forall\, n \ge n_0 .
\end{equation}
 
For all values $n \ge n_0$, $f(n)$ equals $g(n)$ up to constant factors.
\newline
$g(n)$ is an \highlightit{asymptotically tight bound} for $f(n)$.

\subsection{Meaning of previous notations - Insertion sort case}
\textbf{Big O notation}
\newline
describes an \highlightit{upper bound}.
When we use it to \highlightit{bound the worst-case running time} of an 
algorithm, we have a bound on the running time of the algorithm on \highlightit{every input}.
\newline
Thus, the $O(n^2)$ bound on worst-case running time of \highlightit{insertion sort} 
also applies to its running time on every input.
\newline
\textbf{Omega notation} 
\newline
describes a \highlightit{lower bound}.
When we say that the running time (no modifier) of an algorithm is $\Omega(g(n))$, 
we mean that 
\highlightit{no matter what particular input of size $n$ is chosen} for
each value of $n$, the running time on that input is at least a constant 
times $g(n)$, for sufficiently large $n$.
\newline
Equivalently, we are giving a \highlightit{lower bound on the best-case running time} 
of an algorithm. For example, the best-case running time of insertion 
sort is $\Omega(n)$, which implies that the running time of insertion sort 
is $\Omega(n)$.
\newline
The \highlightit{running time of insertion sort therefore belongs to both $\Omega(n)$ 
and $O(n^2)$}, since it falls anywhere between a linear function of $n$ and 
a quadratic function of $n$.
Moreover, these bounds are asymptotically as tight as possible: 
for instance, the running time of insertion sort is not $\Omega(n^2)$, 
since there exists an input for which 
insertion sort runs in $\Theta(n)$ time 
(e.g., when the input is already sorted). 
It is not contradictory, however, 
to say that the worst-case running time of insertion sort is $\Omega(n^2)$, 
since there exists an input that causes the algorithm to take $\Omega(n^2)$ time.
\newline
\textbf{Theta notation}
\newline{}
describes a \highlightit{tight bound}. The $\Theta(n^2)$ notation bound on the worst-case running time of insertion sort, however, 
does not imply a $\Theta(n^2)$ bound 
on the running time of insertion sort on every input. 
For example, when the input is already sorted, 
insertion sort runs in $\Theta(n)$ time.
Technically, 
it is an abuse to say that the running time of insertion sort is $O(n^2)$, 
since for a given $n$, 
the actual running time varies, depending on the particular input of size $n$. 
When we say ``the running time is $O(n^2)$'', 
we mean that 
there is a function $f(n)$ that is $O(n^2)$ 
such that for any value of $n$, 
no matter what particular input of size $n$ is chosen, 
the running time on that input is bounded from above by the value $f(n)$.
Equivalently, we mean that the worst-case running time is $O(n^2)$.

\newpage
\subsection{Running times of well known algorithms}

\begin{table}[H]
    \small
    \begin{tabularx}{\textwidth}{|l|c|c|c|>{\raggedright\arraybackslash}X|}
        \hline
        Algorithm & Best & Average & Worst & Extra space \\
        \hline
        Quicksort & $O(n\log n)$ & $O(n\log n)$ & $O(n^{2})$ &
          $O(\log n)$ avg stack; $O(n)$ worst \\
        \hline
        Mergesort & $O(n\log n)$ & $O(n\log n)$ & $O(n\log n)$ & $O(n)$ \\
        \hline
        Timsort & $O(n)$ & $O(n\log n)$ & $O(n\log n)$ & $O(n)$ \\
        \hline
        Heapsort & $O(n\log n)$ & $O(n\log n)$ & $O(n\log n)$ &
          $O(1)$ aux.\ (in-place) \\
        \hline
        Bubble Sort & $O(n)$ & $O(n^{2})$ & $O(n^{2})$ & $O(1)$ \\
        \hline
        Insertion Sort & $O(n)$ & $O(n^{2})$ & $O(n^{2})$ & $O(1)$ \\
        \hline
        Selection Sort & $O(n^{2})$ & $O(n^{2})$ & $O(n^{2})$ & $O(1)$ \\
        \hline
        Shell Sort & gap-dep.$^{\dagger}$ & gap-dep.$^{\dagger}$ &
          $\le O(n^{2})^{\dagger}$ & $O(1)$ \\
        \hline
        Bucket Sort & $O(n+k)^{\ast}$ & $O(n+k)^{\ast}$ &
          $O(n^{2})^{\ddagger}$ & $O(n+k)$ \\
        \hline
        Radix Sort & $O(d(n+r))$ & $O(d(n+r))$ & $O(d(n+r))$ &
          $O(n+r)$ \\
        \hline
    \end{tabularx}
\end{table}
{\footnotesize
$^{\ast}$Bucket: roughly uniform keys, $k$ buckets; also depends on per-bucket sort.
$^{\ddagger}$Worst when nearly all keys fall in one bucket sorted quadratically.
$^{\dagger}$Shellsort: no universal closed form --- depends on the gap sequence; poor gaps stay $O(n^{2})$.
Radix: $d$ = digits/passes, $r$ = radix (digit range).
}

Verbal description of the previous sorting algorithms:
\begin{itemize}

    \item \href{https://en.wikipedia.org/wiki/Quicksort}{\highlight{Quicksort}} \\
    Efficient general-purpose sorting algorithm. On \highlight{randomized} or typical data it 
    is often \highlight{faster in practice} than mergesort/heapsort thanks to locality and 
    small constants --- that is an empirical claim, not a theorem. 
    Worst-case time is $O(n^2)$ (e.g.\ adversarial pivots); randomized or median-of-three 
    pivots make that rare. \\
    It selects a \highlight{pivot} and \highlight{partitions} the other elements into 
    \highlight{two sub-arrays} (less / greater than the pivot) --- hence 
    \highlight{partition-exchange sort}. Sub-arrays are then sorted recursively. 
    Stack / recursion depth is $O(\log n)$ on balanced partitions and $O(n)$ in the 
    degenerate case.
    
    \item \href{https://en.wikipedia.org/wiki/Merge_sort}{\highlight{Mergesort}} \\
    \highlight{Divide}: split the array into two halves and sort each half recursively 
    until runs of length one. \highlight{Conquer}: \highlight{merge} two already-sorted 
    halves into one sorted run by walking both with two pointers. 
    Stable; guaranteed $O(n\log n)$ time; needs $\Theta(n)$ extra buffer for a standard merge.
        
    \item \href{https://en.wikipedia.org/wiki/Timsort}{\highlight{Timsort}} \\
    A hybrid, \highlight{stable} sort derived from \highlight{mergesort} and 
    \highlight{insertion sort}, designed for real-world partially ordered data. 
    It finds already-ordered \highlight{runs} and merges them under stack invariants. 
    Implemented by Tim Peters (2002) for CPython; it remains the algorithm behind 
    Python's \texttt{list.sort} / \texttt{sorted} (including current 3.x).
    
    \item \href{https://en.wikipedia.org/wiki/Heapsort}{\highlight{Heapsort}} \\
    Comparison-based sort that can be seen as \highlight{selection sort on a heap}: 
    grow a sorted region by repeatedly extracting the max from a max-heap over the 
    unsorted region. Unlike naive selection sort, finding the max is $O(\log n)$ via 
    the heap, not a linear scan. In-place with $O(1)$ \emph{auxiliary} space; not stable.
    
    \item \href{https://en.wikipedia.org/wiki/Bubble_sort}{\highlight{Bubble Sort}} \\
    Repeatedly passes through the list, \highlight{comparing adjacent} elements and 
    swapping if out of order, until a pass makes no swaps (list sorted). Best case $O(n)$ 
    with an early-exit flag; average/worst $O(n^2)$.
    
    \item \href{https://en.wikipedia.org/wiki/Insertion_sort}{\highlight{Insertion Sort}} \\
    Grows a sorted prefix: each step takes the next unsorted element and 
    \highlight{inserts} it into the correct place in the sorted prefix (shifting larger 
    elements). Excellent on nearly sorted data and small $n$; used inside Timsort for short runs.
    
    \item \href{https://en.wikipedia.org/wiki/Selection_sort}{\highlight{Selection Sort}} \\
    Repeatedly finds the \highlight{smallest} remaining unsorted element and swaps it 
    into the next position of the sorted prefix. Always $\Theta(n^2)$ comparisons; 
    few writes; not stable in the usual swap form.
    
    \item \href{https://en.wikipedia.org/wiki/Shellsort}{\highlight{Shell Sort}} \\
    Insertion sort on interleaved subarrays with decreasing \highlight{gap} sizes 
    (e.g.\ $n/2, n/4, \ldots, 1$). Early large gaps move elements far; the last gap-1 
    pass is ordinary insertion sort on a nearly sorted array. 
    Complexity \highlight{depends on the gap sequence}: some sequences give 
    $o(n^2)$ worst case; poor sequences remain $O(n^2)$. No single universal 
    average-case formula --- quote the sequence when you quote a bound.
    
    \item \href{https://en.wikipedia.org/wiki/Bucket_sort}{\highlight{Bucket Sort}} \\
    Scatter keys into $k$ \highlight{buckets} by range, sort each bucket 
    (often insertion sort or recursion), then concatenate. A \highlight{distribution} 
    sort (related to pigeonhole / radix ideas). Best when keys are roughly uniform; 
    if almost all keys land in one bucket and that bucket is sorted quadratically, 
    time falls to $O(n^2)$. Comparisons appear only inside buckets if you use a 
    comparison sort there --- bucket sort itself is not ``just'' a comparison sort.
    
    \item \href{https://en.wikipedia.org/wiki/Radix_sort}{\highlight{Radix Sort}} \\
    Sort integers (or fixed-length keys) \highlight{digit by digit} 
    (LSD or MSD), typically with a stable counting pass per digit. 
    Time $O(d(n+r))$ for $d$ digits and radix $r$ (e.g.\ $r=256$ for bytes). 
    Not a comparison sort; excellent when $d$ is small relative to $\log n$.

\end{itemize}

\subsection{Methods for solving recurrences}
See~\cite{algo-clrs}. Divide-and-conquer running times often satisfy a recurrence
$T(n) = a\, T(n/b) + f(n)$ (plus a base case). Three classic ways to solve them:

\toclessnparagraph{Substitution}
Guess a closed form, then prove it by induction (strengthen the induction hypothesis
if additive lower-order terms get in the way).
Example --- binary search: $T(n) = T(\lfloor n/2\rfloor) + \Theta(1)$.
Guess $T(n) \le c\log n$ for large enough $c$; the inductive step is immediate once
the base case and $c$ absorb the $\Theta(1)$ work. Result: $T(n) = \Theta(\log n)$.

\toclessnparagraph{Recursion tree}
Draw the tree of recursive calls; sum the non-recursive cost $f(\cdot)$ at each level,
then sum over levels (plus leaves if needed).
Example --- mergesort: $T(n) = 2T(n/2) + \Theta(n)$.
Each level costs $\Theta(n)$ in total; there are $\Theta(\log n)$ levels until size-1
leaves --- so $T(n) = \Theta(n\log n)$. Good for building intuition before a formal proof.

\toclessnparagraph{Master theorem}
For $T(n) = a\, T(n/b) + f(n)$ with $a \ge 1$, $b > 1$, compare $f(n)$ to
$n^{\log_b a}$ (let $c_{\mathrm{crit}} = \log_b a$):
\begin{itemize}
    \item \textbf{Case 1:} $f(n) = O\!\bigl(n^{c_{\mathrm{crit}}-\varepsilon}\bigr)$ for some
          $\varepsilon > 0$ $\Rightarrow$ $T(n) = \Theta\!\bigl(n^{c_{\mathrm{crit}}}\bigr)$
          (leaves dominate).
    \item \textbf{Case 2:} $f(n) = \Theta\!\bigl(n^{c_{\mathrm{crit}}}\bigr)$
          $\Rightarrow$ $T(n) = \Theta\!\bigl(n^{c_{\mathrm{crit}}}\log n\bigr)$
          (levels cost about the same; $\Theta(\log n)$ of them).
    \item \textbf{Case 3:} $f(n) = \Omega\!\bigl(n^{c_{\mathrm{crit}}+\varepsilon}\bigr)$
          and a regularity condition on $f$
          $\Rightarrow$ $T(n) = \Theta(f(n))$ (root work dominates).
\end{itemize}
Quick checks: mergesort $a{=}2$, $b{=}2$, $f(n){=}\Theta(n)$ is Case~2
$\Rightarrow \Theta(n\log n)$; binary search $a{=}1$, $b{=}2$, $f(n){=}\Theta(1)$ is
Case~2 with $c_{\mathrm{crit}}{=}0$ $\Rightarrow \Theta(\log n)$.
Not every recurrence fits --- when $a,b$ vary or $f$ is awkward, fall back to
substitution or the tree.

\subsection{Comparison-sort lower bound}
In the algebraic decision-tree / comparison model, any correct comparison-based
sort needs $\Omega(n\log n)$ comparisons in the worst case (there are $n!$ permutations;
each comparison has two outcomes, so the decision tree height is
$\Omega(\log(n!)) = \Omega(n\log n)$ by Stirling).
Hence mergesort and heapsort are \highlight{asymptotically optimal} among comparison sorts;
radix/bucket can beat $n\log n$ only by using the \emph{structure of the keys}, not
comparisons alone.

\section{Space Complexity}
\label{sec:space-complexity}
\highlight{Space complexity} counts memory used as a function of input size $n$.
Distinguish:
\begin{itemize}
    \item \textbf{Input space} --- the array/structure you are given.
    \item \textbf{Auxiliary (extra) space} --- stack frames, buffers, heaps beyond the input.
    \item \textbf{Total space} --- input + auxiliary.
\end{itemize}
Examples: in-place heapsort / insertion sort use $O(1)$ auxiliary space; 
standard mergesort uses $\Theta(n)$ auxiliary for the merge buffer; 
quicksort's auxiliary cost is mainly the \highlight{call stack} 
($O(\log n)$ typical, $O(n)$ worst). 
For interviews, say whether you mean auxiliary or total, and whether recursion counts.

\section{Order Book}
\label{sec:order-book}
An \highlight{order book} stores resting interest: buy and sell orders waiting to trade. 
Exchanges and many trading systems publish incremental updates (add / modify / cancel / trade); 
your process rebuilds a local book and reads \highlight{best bid/offer} (BBO) and depth for 
decisions.
\par\noindent\textbf{Interview tip:} focus on data layout, price-time priority, and hot-path cost --- not exchange 
matching-engine internals.

\subsection{Price Levels and Market Depth}
\label{subsec:order-book-price-levels}
Orders at the \highlight{same price} aggregate into a \highlight{price level}: typically 
price + total quantity (+ optional order count or per-order queue).
\newline
\highlightbf{Market depth} is how many levels you keep on each side (top of book only vs 
$N$ levels vs full book). In market-data vocabulary that lines up roughly with 
\highlightbf{L1} (top of book / BBO), \highlightbf{L2} (depth by price: aggregated quantity 
per level), and \highlightbf{L3} (depth by order: each resting order id). Exact feed labels 
vary by venue; you can also subscribe to L2/L3 and still \highlight{store} only L1 locally. 
Deeper books cost more memory, more update work, and worse cache behavior; many strategies 
only need top-of-book or a small depth window.

\subsection{Bids and Asks}
\label{subsec:order-book-bids-asks}
\begin{itemize}
    \item \highlightbf{Bid} --- buy interest. Best bid = \highlight{highest} buy price.
    \item \highlightbf{Ask} (offer) --- sell interest. Best ask = \highlight{lowest} sell price.
\end{itemize}
A healthy book has best bid $<$ best ask. The difference is the \highlight{spread}; 
midpoint is often $(\mathrm{best\ bid} + \mathrm{best\ ask})/2$. 
Crossed or locked books (bid $\ge$ ask) usually mean a bad feed, race, or auction state --- 
worth detecting in production.

\subsection{Data Structures}
\label{subsec:order-book-data-structures}
See Chapter~\ref{chap:containers} for container APIs. 
This subsection compares order-book-specific representations only.
\newline
You need: fast \highlight{find level by price}, fast \highlight{BBO}, and (often) 
walk levels in price order. Common choices:
\begin{itemize}
    \item \highlightbf{Ordered map} (\texttt{std::map} / flat map): price $\rightarrow$ level. 
          Logarithmic update; natural ordered iteration. Simple; often fine for whiteboard 
          sketches and moderate update rates.
    \item \highlightbf{Hash map + sorted structure}: \texttt{unordered\char`_map} for 
          price $\rightarrow$ level, plus a heap or side tree for best price. Faster random 
          price lookup; more moving parts to keep consistent.
    \item \highlightbf{Price-indexed array / sparse ladder}: if the instrument has a discrete 
          tick grid and a bounded price range, index by tick. Best-case $O(1)$ level access; 
          careful with range and memory.
    \item \highlightbf{Per-level order queue}: if you need individual order ids (modify/cancel 
          by id), each level holds a FIFO list (price-time priority). Aggregate-only books 
          store quantity per price and skip per-order lists.
\end{itemize}
\begin{lstlisting}[language=C++]
#include <cstdint>
#include <map>

struct Level {
    int64_t qty{0};
    uint32_t orders{0};
};

// Bids: highest price first. Asks: lowest price first.
std::map<int64_t, Level, std::greater<int64_t>> bids;
std::map<int64_t, Level, std::less<int64_t>>    asks;
\end{lstlisting}
\par\noindent\textbf{Interview tip:} say \highlight{what} you optimize (BBO vs deep book vs cancel-by-id) before 
picking the structure.

\subsection{Add, Update, and Remove}
\label{subsec:order-book-operations}
Feed handlers typically apply:
\begin{itemize}
    \item \highlightbf{Add} --- new order at a price (create level or increase qty / enqueue order).
    \item \highlightbf{Modify} / replace --- change qty or price (often cancel+add semantics).
    \item \highlightbf{Cancel} / delete --- remove qty or whole order; drop empty levels.
    \item \highlightbf{Trade} / execute --- reduce resting qty when a fill hits the book 
          (depending on whether the feed already netted that into cancel/modify).
\end{itemize}
Maintain invariants: no empty levels left in the maps; BBO caches invalidated or updated when 
the top level changes; order-id index consistent if you track individual orders.
\newline
Sequence numbers / gaps: market-data feeds are often UDP; detect gaps and rebuild from snapshot 
(Chapter~\ref{chap:linuxnetworking}).

\subsection{Best Bid and Offer}
\label{subsec:order-book-bbo}
\highlightbf{BBO} = best bid price/qty and best ask price/qty. Hot strategies read this every 
tick. Keep it as an explicit cached pair updated on each top-of-book change rather than 
re-scanning the book.
\begin{lstlisting}[language=C++]
struct BBO {
    int64_t bid_px{0}, bid_qty{0};
    int64_t ask_px{0}, ask_qty{0};
};

BBO top(const std::map<int64_t, Level, std::greater<int64_t>>& bids,
        const std::map<int64_t, Level, std::less<int64_t>>& asks)
{
    BBO b{};
    if (!bids.empty()) {
        b.bid_px = bids.begin()->first;
        b.bid_qty = bids.begin()->second.qty;
    }
    if (!asks.empty()) {
        b.ask_px = asks.begin()->first;
        b.ask_qty = asks.begin()->second.qty;
    }
    return b;
}
\end{lstlisting}

\subsection{Hot-Path Design}
\label{subsec:order-book-performance}
\begin{itemize}
    \item \highlightbf{No allocations} on the update path --- pre-sized maps/pools, recycled 
          nodes (Section~\ref{sec:avoiding-allocations}).
    \item \highlightbf{Contiguous levels} and stable layout beat pointer-chasing lists when 
          depth walks matter (Chapter~\ref{chap:cacheperformance}).
    \item \highlightbf{Single writer} for a given instrument book when possible (SPSC into the 
          book thread --- Section~\ref{sec:spsc-ring-buffer}); avoid locking inside the decode 
          loop.
    \item \highlightbf{NUMA / pinning}: build the book on the core that consumes it 
          (Chapters~\ref{chap:numa},~\ref{chap:lowlatency}).
    \item Measure update latency and cache misses under realistic burst rates; asymptotic 
          map cost is not enough if allocation jitter dominates.
\end{itemize}
